\section{Stieltjes collisions at different powers of the small parameter}
\label{mc:section}

The unequal-rate collision theorem in the latest gap-six report
\cite[Section~4 and research question~10]{GapSix} assumes
$a=\lambda\epsilon$ and $b=\mu\epsilon$ with fixed positive rates.
Its proof explicitly leaves different powers of $\epsilon$ open.
Here we resolve that question for every pair of positive monomial
scales and every fixed pair of Stieltjes indices and argument
derivatives. The additional constants occur at explicitly classified
rational ratios of the exponents. In particular, the product of the
two singular parts can now contribute a constant.

The local subtraction argument is a classical asymptotic method;
related generalized Stieltjes-transform methods appear in
\cite{TicaGalapon}. The contribution here is the complete resonance
classification and its evaluated contact formula for this family,
rather than a claim to invent finite-part asymptotics.

\subsection{A compact identity before the general theorem}

Write $\gamma$ for Euler's constant, and put
\begin{equation}\label{mc:D}
 D(a,b)=\frac{\log(b/a)}{b-a}
       =\int_0^\infty\frac{dx}{(x+a)(x+b)},\qquad a,b>0.
\end{equation}
The value at $a=b$ is $1/a$. A derivative in the first argument is
written $D_a$; explicitly,
\begin{equation}\label{mc:Da}
 D_a(a,b)=-\frac1{a(b-a)}+
                    \frac{\log(b/a)}{(b-a)^2}.
\end{equation}

\begin{theorem}[A two-scale ordered digamma identity]\label{mc:ordered-limit}
Fix $\lambda,\mu,\alpha,\beta>0$, and set
$a=\lambda\epsilon^\alpha$, $b=\mu\epsilon^\beta$.
Then, as $\epsilon\downarrow0$,
\begin{equation}\label{mc:ordered-uniform}
 \boxed{\displaystyle
 \int_0^1\psi'(x+a)\psi(x+b)\,dx
 =D_a(a,b)+\frac{\psi(1+b)}a
                   +\zeta(2)\log\frac ba+o(1).}
\end{equation}
Every term on the left is an ordinary convergent integral. The
subtraction on the right is exact to a remainder tending to zero
along every indicated monomial path.
\end{theorem}

\begin{proof}
Let $h(x)=\psi(1+x)$, so that $\psi(x)=-1/x+h(x)$.
Expand the product and integrate the term
$h(x+b)/(x+a)^2$ by parts. The difference between the two sides of
\eqref{mc:ordered-uniform}, without $o(1)$, is
\begin{align}\label{mc:ordered-remainder}
 &-D_a(1+a,1+b)-\frac{h(1+b)}{1+a}
     +\zeta(2)\log\frac{1+a}{1+b}\notag\\
 &\quad+\int_0^1\left\{
       \frac{h'(x+b)-\zeta(2)}{x+a}
      -\frac{h'(x+a)-\zeta(2)}{x+b}\right\}dx
      +\int_0^1h'(x+a)h(x+b)\,dx.
\end{align}
In the middle integral replace $h'(x+b)$ and $h'(x+a)$ by
$h'(x)$. The two errors are
$O(b(1+|\log a|)+a(1+|\log b|))$.
Since $h'(x)-\zeta(2)=O(x)$, the remaining difference tends to
zero by dominated convergence. The first, second, and last terms
in \eqref{mc:ordered-remainder} tend respectively to
\[
 \frac12,\qquad \gamma-1,\qquad
 \frac12\bigl(\psi(2)^2-\psi(1)^2\bigr)=\frac12-\gamma.
\]
They cancel. A positive power of $\epsilon$ times any fixed power
of $|\log\epsilon|$ tends to zero, proving the statement.
\end{proof}

For a function with a power--logarithm expansion, let
$\operatorname{CT}_{\epsilon}$ denote the coefficient of
$\epsilon^0(\log\epsilon)^0$. It does not mean the ordinary limit
of a divergent integral. Let $\mathbf1_{\mathcal E}$ denote the
indicator of the stated condition.

\begin{corollary}[An explicit resonance law]\label{mc:ordered-CT}
Suppose $q=\alpha/\beta>1$, and set $j=2/(q-1)$.
For $I(\epsilon)$ equal to the integral in
\eqref{mc:ordered-uniform},
\begin{align}\label{mc:ordered-constant}
 \operatorname{CT}_{\epsilon}I
 ={}&\zeta(2)\log\frac\mu\lambda
  +\mathbf1_{q\in\mathbb Z_{\ge2}}
       (-1)^{q+1}\frac{\mu^q}{\lambda}\zeta(q+1)\notag\\
 &+\mathbf1_{j\in\mathbb Z_{\ge1}}
       \frac{\lambda^j}{\mu^{j+2}}
                  \left((j+1)\log\frac\mu\lambda-1\right).
\end{align}
A term whose indicator is zero is omitted, including its
integer-indexed notation. In particular, at $\lambda=\mu=1$,
\begin{equation}\label{mc:ordered-examples}
 \operatorname{CT}_{\epsilon}I=
 \begin{cases}
  -1-\zeta(3),&q=2,\\
  -1+\zeta(4),&q=3,\\
  (-1)^{q+1}\zeta(q+1),&q\in\mathbb Z,\ q\ge4,\\
  -1,&q=1+2/j,\ j\ge3,\\
  0,&q\notin\mathbb Z_{\ge2}\ \text{and }2/(q-1)\notin\mathbb Z_{\ge1}.
 \end{cases}
\end{equation}
\end{corollary}

\begin{proof}
Use the convergent expansion
$\psi(1+b)=-\gamma+\sum_{\ell\ge1}(-1)^{\ell+1}
\zeta(\ell+1)b^\ell$ for $|b|<1$.
In $\psi(1+b)/a$, exponent zero requires $\ell\beta=\alpha$.
If $\rho=a/b$, the other singular term has the convergent expansion
\begin{equation}\label{mc:Da-series}
 D_a(a,b)=-\frac1{ab}
 +\frac1{b^2}\sum_{k\ge0}\rho^k
                   \left((k+1)\log\frac ba-1\right).
\end{equation}
Here exponent zero requires $k(\alpha-\beta)=2\beta$.
Extracting the constant in $\log\epsilon$ proves the result.
\end{proof}

These are new exact values of specified asymptotic coefficients.
They also give ordinary renormalized limits by subtracting the
finitely many divergent powers and logarithms in
\eqref{mc:Da-series} and the displayed digamma series. For example,
if $L=\log\epsilon$, the first two cases are
\begin{align}\label{mc:q2limit}
 \lim_{\epsilon\downarrow0}\bigg[
  \int_0^1\psi'(x+\epsilon^2)\psi(x+\epsilon)dx
   +\epsilon^{-3}+(1+\gamma+L)\epsilon^{-2}\notag\\
 \hspace{28mm}
   +(1+2L-\zeta(2))\epsilon^{-1}
   +(3+\zeta(2))L\bigg]
 &=-1-\zeta(3),\\[3pt]
 \lim_{\epsilon\downarrow0}\bigg[
  \int_0^1\psi'(x+\epsilon^3)\psi(x+\epsilon)dx
   +\epsilon^{-4}+\gamma\epsilon^{-3}
   +(1+2L-\zeta(2))\epsilon^{-2}\notag\\
 \hspace{28mm}
   +\zeta(3)\epsilon^{-1}
   +(4+2\zeta(2))L\bigg]
 &=-1+\zeta(4).
 \label{mc:q3limit}
\end{align}
The convergence is established analytically by the theorem;
the numerical diagnostics in the package test the signs and
the independent expansions.

\subsection{Definitions for arbitrary Stieltjes indices}

We fix the standard convention
\begin{equation}\label{mc:stieltjes}
 Z(u,x)=\zeta(1+u,x)-\frac1u
     =\sum_{m\ge0}\frac{(-1)^m}{m!}\gamma_m(x)u^m,
 \qquad \gamma_0(x)=-\psi(x).
\end{equation}
Argument derivatives are indicated by a parenthesized superscript.
Put
\begin{equation}\label{mc:notation}
 \begin{gathered}
 P_r(u)=(-1)^r(1+u)_r,\qquad
 A_r(u)=\frac{(1+u)_r}{r!},\\
 Z_r(u,x)=\partial_x^r Z(u,x),\qquad
 g_h(u)=\left.\partial_x^h Z(u,x)\right|_{x=1}.
 \end{gathered}
\end{equation}
These are holomorphic near $u=0$. The Hurwitz recurrence
\cite[\S25.11]{DLMF} gives the exact splitting
\begin{equation}\label{mc:split}
       Z_r(u,x)=P_r(u)x^{-r-1-u}+Z_r(u,1+x).
\end{equation}
In particular $g_0(u)=\zeta(1+u)-1/u$ and, for $h\ge1$,
\begin{equation}\label{mc:g-explicit}
        g_h(u)=(-1)^h(1+u)_h\zeta(1+u+h).
\end{equation}
Thus every derivative of $g_h$ is a finite combination of
ordinary zeta derivatives and generalized harmonic numbers;
for $h=0$ its coefficients are the ordinary Stieltjes constants.

The reference finite part is
\begin{equation}\label{mc:FP}
 F_{m,n}^{r,t}=\operatorname{CT}_{\delta\downarrow0}
  \int_\delta^1\gamma_m^{(r)}(x)\gamma_n^{(t)}(x)\,dx.
\end{equation}
It uses the coordinate $\delta$ exactly as displayed. An explicit
holomorphic germ $\mathcal F_{r,t}(u,v)$ for these coefficients
can be defined entirely by an ordinary convergent integral.
With $R=r+t+1$, set
\begin{align}\label{mc:FP-remainder}
 \mathcal R(u,v;x)={}&Z_r(u,x)Z_t(v,x)
       -P_r(u)P_t(v)x^{-R-1-u-v}\notag\\
 &-P_r(u)\sum_{k=0}^r\frac{g_{t+k}(v)}{k!}x^{k-r-1-u}
  -P_t(v)\sum_{k=0}^t\frac{g_{r+k}(u)}{k!}x^{k-t-1-v},\\
 \mathcal F_{r,t}(u,v)={}&\int_0^1\mathcal R(u,v;x)dx
  -\frac{P_r(u)P_t(v)}{R+u+v}\notag\\
 &-P_r(u)\sum_{k=0}^{r-1}\frac{g_{t+k}(v)}{k!(r-k+u)}
  -P_t(v)\sum_{k=0}^{t-1}\frac{g_{r+k}(u)}{k!(t-k+v)}.
 \label{mc:FP-germ}
\end{align}
Empty sums vanish. Then
$F_{m,n}^{r,t}=(-1)^{m+n}\partial_u^m\partial_v^n
\mathcal F_{r,t}(0,0)$.
The critical powers $x^{-1}\log^h x$ contribute zero to the
constant on $(0,1)$; this explains their omission from the two
rational sums in \eqref{mc:FP-germ}. This construction agrees with
the inherited convention \cite{GapSix,LatticeIncoming}.

For example, those sources give
\begin{equation}\label{mc:FP-values}
 \begin{gathered}
 F_{0,0}^{0,0}=2\gamma_1-2\zeta(2),\quad
 F_{0,0}^{1,0}=\zeta(2),\quad
 F_{0,0}^{1,1}=2\zeta(3)+4\zeta'(3),\\
 F_{1,0}^{0,0}=\tfrac32\gamma_2+2\gamma\zeta(2)-\zeta(3).
 \end{gathered}
\end{equation}
These baseline evaluations are credited to the previous work.
The new result evaluates every change caused by the different scales.

\subsection{The complete constant formula}

Assume henceforth that $\alpha>\beta>0$ and define
\begin{equation}\label{mc:scales}
 a=\lambda\epsilon^\alpha,\quad b=\mu\epsilon^\beta,
 \quad d=\alpha-\beta,\quad \tau=\lambda/\mu,
 \quad R=r+t+1.
\end{equation}
The opposite ordering follows by exchanging the factors.
Introduce the finite set
\begin{equation}\label{mc:resonance-set}
 \mathscr L_r(\alpha,\beta)
 =\{(k,\ell):0\le k<r,\ \ell\in\mathbb Z_{\ge1},\
                      \beta\ell=\alpha(r-k)\}.
\end{equation}
It has at most $r$ members, including at rational ratios with
large numerator or denominator. Define the following removable germ:
\begin{align}\label{mc:mixed-contact}
 \mathcal C_{r,t}(u,v;\lambda,\mu;\alpha,\beta)
 ={}&(-1)^r g_{r+t}(v)
                  \frac{\lambda^{-u}-A_r(u)}u\notag\\
 &+(-1)^r\sum_{(k,\ell)\in\mathscr L_r(\alpha,\beta)}
   \frac{\mu^\ell\lambda^{k-r-u}}{\ell!}
                 (1+u)_{r-k-1}\,g_{t+k+\ell}(v).
\end{align}
The first quotient is removable since $A_r(0)=1$.

There is one possible singular--singular contact. If
$j=\beta R/d$ is a positive integer, set
\begin{align}\label{mc:singular-contact}
 \mathcal H_j(u,v)={}&
 \frac{(-1)^t\mu^{-R-v}\tau^j}{u\Gamma(1+v)}\notag\\
 &\quad\times\left\{
 \frac{\lambda^{-u}\Gamma(R+j+v)}{(1-u)_j}
 -\frac{\mu^{-u}\Gamma(1-u)\Gamma(R+j+u+v)}{j!}
                 \right\}.
\end{align}
The numerator vanishes identically at $u=0$, so this is a
holomorphic two-variable germ. Define $\mathcal H_*=\mathcal H_j$
when $\beta R/d=j\in\mathbb Z_{\ge1}$, and $\mathcal H_*=0$
otherwise. The quotient must be removed before taking derivatives.

\begin{theorem}[All-index monomial collision constants]\label{mc:main}
For every $m,n,r,t\in\mathbb Z_{\ge0}$ and all positive parameters
in \eqref{mc:scales}, the integral
\[
 I_{m,n}^{r,t}(\epsilon)
   =\int_0^1\gamma_m^{(r)}(x+a)
                   \gamma_n^{(t)}(x+b)\,dx
\]
has a power--logarithm asymptotic expansion through its constant.
Its exact constant is
\begin{align}\label{mc:main-constant}
 \boxed{\begin{aligned}
 \operatorname{CT}_{\epsilon}I_{m,n}^{r,t}
  =F_{m,n}^{r,t}+(-1)^{m+n}\partial_u^m\partial_v^n
    \big[&\mathcal C_{r,t}(u,v;\lambda,\mu;\alpha,\beta)\\
       +&\mathcal C_{t,r}(v,u;\mu,\lambda;\beta,\alpha)
         +\mathcal H_*(u,v)\big]_{u=v=0}.
 \end{aligned}}
\end{align}
Every correction in this formula is a finite expression in
logarithms, ordinary zeta values and derivatives, Stieltjes constants
at one, and generalized harmonic numbers. There is no new
unevaluated spectral constant in the correction.
\end{theorem}

There are precisely three possible mechanisms for extra constants:
\begin{equation}\label{mc:three-mechanisms}
 \beta\ell=\alpha(r-k),\qquad
 \alpha\ell=\beta(t-k),\qquad
 j(\alpha-\beta)=\beta(r+t+1).
\end{equation}
The two first equations have the finite ranges specified above;
the last has at most one admissible integer $j$. If $\alpha/\beta$
is irrational, all three resonance contributions vanish. Only the
first, critical quotient in each $\mathcal C$ remains. These are
selection rules for an asymptotic coefficient, not statements of
discontinuity of the original positive-shift integral.

\subsection{A two-shift expansion with analytic coefficients retained}

We prove the theorem in two steps. Keeping the analytic coefficients
at their actual small shifts is the useful replacement for a
fixed-ratio Taylor truncation. Define
\begin{equation}\label{mc:kernel}
 \mathcal B_{r,t}(u,v;a,b)
  =P_r(u)P_t(v)\int_0^\infty
             (x+a)^{-r-1-u}(x+b)^{-t-1-v}\,dx.
\end{equation}
This is an ordinary absolutely convergent integral for $u,v$ near
zero. It also has the elementary incomplete-beta representation
used below.

\begin{lemma}[Uniform germ along monomial paths]\label{mc:uniform-germ}
For fixed exponents and rates as in \eqref{mc:scales}, locally
uniformly on some complex neighborhood of $(u,v)=(0,0)$,
\begin{align}\label{mc:uniform-expansion}
 &\int_0^1 Z_r(u,x+a)Z_t(v,x+b)dx
    =\mathcal F_{r,t}(u,v)+\mathcal B_{r,t}(u,v;a,b)\notag\\
 &\quad+(-1)^r\sum_{k=0}^{r-1}
     (1+u)_{r-k-1}a^{k-r-u}Z_{t+k}(v,1+b)
   +(-1)^r g_{r+t}(v)\frac{a^{-u}-A_r(u)}u\notag\\
 &\quad+(-1)^t\sum_{k=0}^{t-1}
     (1+v)_{t-k-1}b^{k-t-v}Z_{r+k}(u,1+a)
   +(-1)^t g_{r+t}(u)\frac{b^{-v}-A_t(v)}v+o(1).
\end{align}
All fixed spectral derivatives of the remainder tend to zero.
\end{lemma}

\begin{proof}
Use \eqref{mc:split}. The singular--singular product on $(0,1)$
equals \eqref{mc:kernel} minus its tail on $(1,\infty)$; the latter
tends uniformly to $P_r(u)P_t(v)/(R+u+v)$.
For the first mixed product, use Taylor's theorem in $x$ at the
point $1+b$:
\[
 Z_t(v,1+b+x)=\sum_{k=0}^r
          \frac{Z_{t+k}(v,1+b)}{k!}x^k+O(x^{r+1}).
\]
The remainder has a bound uniform for small $b$ and for $v$ in a
fixed disk. After multiplication by the singular factor it is
bounded by $Cx^{-\eta}$, with $\eta<1$, independent of the
relative sizes of $a,b$. Dominated convergence therefore applies.
For $k<r$ the beta integral gives
\begin{align*}
 P_r(u)\int_0^1x^k(x+a)^{-r-1-u}dx
 ={}&(-1)^rk!(1+u)_{r-k-1}a^{k-r-u}\\
 &-\frac{P_r(u)}{r-k+u}+o(1).
\end{align*}
For $k=r$ its removable version is
\[
 P_r(u)\int_0^1x^r(x+a)^{-r-1-u}dx
   =(-1)^rr!\frac{a^{-u}-A_r(u)}u+o(1).
\]
One can first take $\Re u>0$, subtract the tail, and continue;
the divided difference is uniformly holomorphic near zero.
Replacing the coefficient of the critical term at $1+b$ by its
value at one costs $O(b\,a^{-\eta}(1+|\log a|))$ on a small disk.
Choose $\eta>0$ so small that $\beta-\alpha\eta>0$; the error
then tends to zero. For the exchanged term also impose
$\alpha-\beta\eta>0$. The upper endpoints and the remaining
integrals give exactly \eqref{mc:FP-germ}. This proves local
uniformity. Cauchy's integral formula on a slightly larger disk
justifies every prescribed spectral derivative.
\end{proof}

The Taylor expansion
\[
 Z_{t+k}(v,1+b)=\sum_{\ell\ge0}
                       g_{t+k+\ell}(v)\frac{b^\ell}{\ell!}
\]
now shows that a noncritical mixed term has base exponent
$\alpha(k-r)+\beta\ell$. Taking it far enough to have a positive
remainder exponent is always a finite task. Exponent zero gives
$\mathscr L_r(\alpha,\beta)$; the critical quotient contributes
its constant in $\log\epsilon$ by replacing $a^{-u}$ with
$\lambda^{-u}$. This proves exactly the two $\mathcal C$ terms.

\subsection{Pairing the singular branches before differentiation}

Put $A=r+1+u$, $B=t+1+v$, $C=A+B-1=R+u+v$ and
$\rho=a/b$. For $0<a<b$, the substitution
$y=(x+a)/(x+b)$ gives
\begin{equation}\label{mc:incomplete-beta}
 \frac{\mathcal B_{r,t}}{P_r(u)P_t(v)}
  =(b-a)^{-C}\left\{
     \frac{\Gamma(1-A)\Gamma(C)}{\Gamma(B)}
      -\int_0^\rho y^{-A}(1-y)^{C-1}\,dy\right\}.
\end{equation}
At generic nonintegral $A$ the last integral is understood by its
analytic continuation; the original integral in
\eqref{mc:kernel} fixes the branch and the removable values.
Expansion of this incomplete beta function, or the usual
hypergeometric connection formula \cite[\S15.8]{DLMF}, gives
two normally convergent series for $0<\rho<1$:
\begin{align}\label{mc:two-branches}
 \frac{\mathcal B_{r,t}}{P_r(u)P_t(v)}
 ={}&b^{-C}\frac{\Gamma(1-A)}{\Gamma(B)}
             \sum_{j\ge0}\frac{\Gamma(C+j)}{j!}\rho^j\notag\\
 &+a^{1-A}b^{-B}
      \sum_{k\ge0}(-1)^k(B)_k
                         \frac{\Gamma(A-1-k)}{\Gamma(A)}\rho^k.
\end{align}
The individual series have poles at $u=0$. Their sum is
holomorphic, and the poles must be paired at equal base exponents.

The $k<r$ terms of the second branch are holomorphic and have
base exponents
$-\alpha r-\beta(t+1)+k(\alpha-\beta)<0$.
For each $j\ge0$, pair the first-branch term indexed by $j$ with
the second-branch term indexed by $k=r+j$. Their common base
exponent is
\begin{equation}\label{mc:paired-exponent}
                    E_j=-\beta R+j(\alpha-\beta).
\end{equation}
After the Gamma reflection and recurrence identities, the paired
term is $\epsilon^{E_j}\mathcal H_j(u,v;L)$, where $L=\log\epsilon$
and
\begin{align}\label{mc:H-log}
 \mathcal H_j(u,v;L)={}&
 \frac{(-1)^t\mu^{-R-v}\tau^j\epsilon^{-\beta v}}
                   {u\Gamma(1+v)}\notag\\
 &\times\left\{
 \frac{\lambda^{-u}\epsilon^{-\alpha u}\Gamma(R+j+v)}{(1-u)_j}
 -\frac{\mu^{-u}\epsilon^{-\beta u}
              \Gamma(1-u)\Gamma(R+j+u+v)}{j!}\right\}.
\end{align}
For each fixed $L$, this is holomorphic in $(u,v)$.
Its constant in $L$ after any fixed spectral differentiation is
obtained by setting $L=0$ before differentiating; the result is
\eqref{mc:singular-contact}.

Choose a finite integer $J$ with $(J+1)(\alpha-\beta)>\beta R$,
then choose the spectral radius $\eta>0$ small enough that also
$(J+1)(\alpha-\beta)-\beta R-(\alpha+\beta)\eta>0$.
On this sufficiently small spectral disk the tail $j>J$ in the
paired series tends uniformly to zero: the coefficient growth
is at most polynomial in $j$, while $\rho^j$ is geometric.
The divided difference in $u$ introduces only a fixed logarithmic
factor, which can be absorbed by shrinking the spectral disk.
Thus only $E_j=0$ can contribute a constant, and this is exactly
the condition defining $\mathcal H_*$. Together with the preceding
mixed-term analysis this proves Theorem~\ref{mc:main} and also
provides all divergent terms through a finite truncation of
\eqref{mc:uniform-expansion} and \eqref{mc:two-branches}.

\subsection{Harmonic and zeta evaluations of the new contacts}

At spectral degree zero, the singular contact is particularly
simple. With $H_0=0$,
\begin{equation}\label{mc:H00}
 \boxed{\displaystyle
 \mathcal H_j(0,0)=(-1)^t\frac{(R+j-1)!}{j!}
       \frac{\lambda^j}{\mu^{R+j}}
             \left(\log\frac\mu\lambda+H_j-H_{R+j-1}\right).}
\end{equation}
This follows by differentiating the numerator of
\eqref{mc:singular-contact} once in $u$.
The derivative of $\log(1-u)_j^{-1}$ is $H_j$ at zero;
the derivative of
$\log\Gamma(1-u)+\log\Gamma(R+j+u)$ is $H_{R+j-1}$.
Formula \eqref{mc:H00} recovers the second resonance family in
Corollary~\ref{mc:ordered-CT}.

For arbitrary spectral degrees, the finite Taylor computations use
\begin{align}\label{mc:harmonic-jets}
 \log\frac{j!}{(1-u)_j}&=\sum_{h\ge1}\frac{H_j^{(h)}}h u^h,\notag\\
 \log\Gamma(N+u)-\log\Gamma(N)
 &=\psi(N)u+
       \sum_{h\ge2}\frac{(-1)^h}{h}
                 \bigl(\zeta(h)-H_{N-1}^{(h)}\bigr)u^h.
\end{align}
Truncate these convergent expansions at the required degree and
exponentiate, equivalently using complete Bell polynomials.
This proves the stated finite constant alphabet constructively.

For completeness, the singular contact can be evaluated without
leaving even a Gamma derivative in the answer. Put $M=R+j$ and
$D=m+n+1$, and form two finite polynomials:
\begin{align}\label{mc:finite-logarithms}
 Q_D(u,v)={}&-u\log\lambda-v\log\mu
       +\sum_{h=1}^D\frac{H_j^{(h)}}h u^h
       +\sum_{h=1}^D\frac{(-1)^{h+1}H_{M-1}^{(h)}}h v^h,
       \notag\\
 T_D(u,v)={}&-(u+v)\log\mu
       +\sum_{h=1}^D\frac{(-1)^{h+1}H_{M-1}^{(h)}}h(u+v)^h
       \notag\\
 &+\sum_{h=2}^D\frac{\zeta(h)}h
          \left[u^h+(-1)^h\bigl((u+v)^h-v^h\bigr)\right].
\end{align}
Define the finite exponential polynomial
$\mathcal E_D(X)=\sum_{k=0}^D X^k/k!$.
Then the evaluated contact is
\begin{equation}\label{mc:finite-jet}
\begin{split}
 &(-1)^{m+n}\partial_u^m\partial_v^n\mathcal H_j(0,0)\\
 &\quad=
 \frac{(-1)^{m+n+t}m!n!(M-1)!}{j!}
     \frac{\lambda^j}{\mu^M}
 [u^{m+1}v^n]\left\{\mathcal E_D(Q_D)-\mathcal E_D(T_D)\right\}.
\end{split}
\end{equation}
The notation $[u^av^b]$ extracts a coefficient from an ordinary
finite polynomial. To prove the formula, divide the two Gamma products
in \eqref{mc:singular-contact} by their common constant $(M-1)!/j!$.
Their logarithms, including the amplitude factors, have Taylor
polynomials $Q_D$ and $T_D$. Euler's constant cancels in each.
Both polynomials have zero constant term, so terms beyond the $D$th
power of their exponential cannot affect the requested coefficient.
Division by the removable $u$ is exactly extraction at $u^{m+1}$.
This gives a finite expression in harmonic numbers and ordinary zeta
values at every specified pair of Stieltjes indices. The two mixed
contacts are obtained just as finitely by multiplying the polynomials
$(1+u)_k$, the exponential of $-u\log\lambda$, and the zeta jets
in \eqref{mc:g-explicit}.

There is an especially transparent first-Stieltjes example. Take
$r=t=0$, so $R=1$. A singular contact occurs exactly when
$\alpha/\beta=1+1/j$, $j\ge1$. Direct expansion of
\eqref{mc:singular-contact} yields
\begin{align}\label{mc:H10}
 -\partial_u\mathcal H_j(0,0)
 =\frac{\lambda^j}{\mu^{j+1}}\bigg[&\zeta(2)-H_j^{(2)}\notag\\
 &+\frac{(H_j-\log\mu)^2-(H_j-\log\lambda)^2}{2}\bigg].
\end{align}
Consequently, for $a=\epsilon^{j+1}$ and $b=\epsilon^j$,
\begin{equation}\label{mc:stieltjes-tail-contact}
 \boxed{\displaystyle
 \operatorname{CT}_{\epsilon}\int_0^1
          \gamma_1(x+\epsilon^{j+1})\gamma_0(x+\epsilon^j)dx
   =\frac32\gamma_2+2\gamma\zeta(2)-\zeta(3)
                         +\zeta(2)-H_j^{(2)}.}
\end{equation}
The new term is the ordinary trigamma tail
$\zeta(2)-H_j^{(2)}=\psi'(j+1)$.
Thus all positive integer trigamma tails appear as exact
Stieltjes collision contacts.

This family also has a classical polylogarithmic generator:
for $|z|<1$,
\begin{equation}\label{mc:polylog-generator}
 \sum_{j\ge1}\bigl(\zeta(2)-H_j^{(2)}\bigr)z^j
       =\frac{\zeta(2)z-\operatorname{Li}_2(z)}{1-z}.
\end{equation}
Its ordinary primitive can be normalized at zero using
\begin{equation}\label{mc:polylog-primitive}
 \int_0^z\frac1t
       \sum_{j\ge1}\bigl(\zeta(2)-H_j^{(2)}\bigr)t^jdt
   =2\operatorname{Li}_{2,1}(z)-\operatorname{Li}_3(z)
      +(\operatorname{Li}_2(z)-\zeta(2))\log(1-z).
\end{equation}
Here $\operatorname{Li}_{2,1}(z)=
\sum_{n>k\ge1}z^n/(n^2k)$, on its analytic branch at zero.
The generating identity is an elementary interchange of absolutely
convergent sums. The primitive identity is checked by differentiation
using $\frac{d}{dz}\operatorname{Li}_{2,1}(z)
=\operatorname{Li}_{1,1}(z)/z=\log^2(1-z)/(2z)$ together with
the corresponding product derivatives; an independent quadrature check
is recorded with the computational artifacts. In fact classical
polylogarithms suffice: for $0<z<1$, the same primitive equals
\begin{equation}\label{mc:classical-primitive}
 2\zeta(3)-\operatorname{Li}_3(z)-2\operatorname{Li}_3(1-z)
                +\log(1-z)\operatorname{Li}_2(1-z).
\end{equation}
To verify this expression directly, differentiate and use Euler's
dilogarithm reflection identity; its limit at zero is zero. At
$z=1$ it gives the endpoint identity
$\sum_{j\ge1}(\zeta(2)-H_j^{(2)})/j=\zeta(3)$.
The elementary generating and primitive identities are classical
Euler-sum consequences; their role here is to evaluate the newly
identified contact family, not to claim those consequences as
independently new polylogarithm identities.

\subsection{What the different-scale theorem changes}

The fixed-rate source remains correct with its stated hypotheses.
It must not be extrapolated by substituting an
$\epsilon$-dependent rate in a remainder proved only for fixed rates.
Two changes are essential: high Taylor coefficients at one can
survive when $b^\ell/a^{r-k}$ has exponent zero, and the two
singular Frobenius branches can jointly have exponent zero.
Both effects are present in \eqref{mc:ordered-examples}; the
first-index law \eqref{mc:stieltjes-tail-contact} isolates the
second effect without argument derivatives.

The theorem does not assert uniformity as the exponent ratio
approaches a resonance while $\epsilon$ also varies. In that
transition, a factor $\epsilon^{E_j}$ must be retained rather than
replaced by its constant coefficient. Formula \eqref{mc:H-log}
provides the exact local term for such a further analysis.
